\sgaStableTarget{sga6.e06.sec.4}\section*{4. Computation of the \(K^\bullet\) of a flag bundle}

\sgaStableTarget{sga6.e06.sec.4.1}\subsection*{4.1}
Let \(S\) be a scheme and let \(E\) be an \(\mathcal O_S\)-module locally free of rank \(n\). Let
\[
  \pi=(p_1,\ldots,p_k),\qquad p_i\geq 1,
  \qquad \sum_{i=1}^k p_i=n.
\]
The flag functor of type \(\pi\) of \(E\) is the functor on the category of \(S\)-schemes, with values in sets, which to an \(S\)-scheme \(f:X\to S\) associates the set of chains of quotients of \(f^*E\)
\[
  (L_{k-1},\ldots,L_1),
\]
where, for each \(i\), \(L_i\) is a quotient of \(L_{i+1}\), locally free of rank \(p_1+\cdots+p_i\). This functor is representable by an \(S\)-scheme \(D_\pi(E)\), endowed with a canonical sequence of quotients of the inverse image of \(E\), locally free of ranks \(p_1,p_1+p_2,\ldots\). If \(\pi=(1,n-1)\), then \(D_\pi(E)=\mathbb P(E)\); if \(\pi=(r,n-r)\), then \(D_\pi(E)=G_{r,n}(E)\), the Grassmannian of type \((r,n-r)\) of \(E\). If all \(p_i\)'s are equal to \(1\), one obtains the usual flag bundle of \(E\), denoted \(D(E)\).

\sgaStableTarget{sga6.e06.lem.4-2}\textbf{Lemma 4.2.} Let \(S\) be a scheme, let \(E\) be an \(\mathcal O_S\)-module locally free of rank \(n\), and let
\[
  \pi=(p_1,\ldots,p_i,\ldots,p_k),\qquad p_1+\cdots+p_k=n.
\]
Let \(p',p''\geq 1\) be such that \(p'+p''=p_i\), and put
\[
  \pi'=(p_1,\ldots,p_{i-1},p',p'',p_{i+1},\ldots,p_k),
  \qquad \pi''=(p',p'').
\]
Let \(L_1,\ldots,L_{k-1}\) be the canonical quotients of the inverse image of \(E\) on \(D_\pi(E)\), and let \(F\) be the kernel of \(L_i\to L_{i-1}\). Then there is a canonical \(S\)-isomorphism
\[
  D_{\pi'}(E)\simeq D_{\pi''}(F).
\]
The proof, immediate from the functorial definition of flag schemes, is left to the reader.

\sgaStableTarget{sga6.e06.lem.4-3}\textbf{Lemma 4.3.} Let \(A\) be a commutative ring with unit, let \(n>0\), and let \((c_i)_{1\leq i\leq n}\) be elements of \(A\). Let \(C\) be the \(A\)-algebra generated by \(\xi_1,\ldots,\xi_n\), subject to the relations
\[
  \sigma_i(\xi_1,\ldots,\xi_n)=c_i,
  \qquad i=1,\ldots,n,
\]
where \(\sigma_i\) denotes the elementary symmetric functions. Let \(p_1,\ldots,p_k\) be integers \(\geq 1\) whose sum is \(n\). For \(1\leq j\leq k\) put
\[
  \xi^{(j)}=(\xi^{(j)}_1,\ldots,\xi^{(j)}_{p_j}),
  \qquad
  \xi_i^{(j)}=\xi_{\sum_{\alpha<j}p_\alpha+i}.
\]
Put
\[
  c_i^{(j)}=\sigma_i(\xi^{(j)}_1,\ldots,\xi^{(j)}_{p_j}),
  \qquad 1\leq j\leq k,
  \quad 1\leq i\leq p_j,
\]
and let \(B\) be the subalgebra of \(C\) generated by the \(c_i^{(j)}\)'s. Then:

\begin{enumerate}[label=\textup{\roman*)}]
\item The elements
\[
  (\xi^{(1)})^{\alpha(1)}\cdots(\xi^{(k)})^{\alpha(k)},
  \qquad 0\leq \alpha_i^{(j)}\leq p_j-i,
\]
form a basis of \(C\), considered as a \(B\)-module.
\item The generators \(c_i^{(j)}\) of the \(A\)-algebra \(B\) satisfy the \(n\) isobaric relations deduced from
\[
  1+\sum_{i=1}^n c_i
  =
  \prod_{j=1}^{k}\left(1+\sum_{i=1}^{p_j} c_i^{(j)}\right),
\]
where \(c_i\) and \(c_i^{(j)}\) are assigned weight \(i\); these \(n\) relations generate the ideal of relations among the \(c_i^{(j)}\)'s.
\end{enumerate}

\sgaStableTarget{sga6.e06.cor.4-4}\textbf{Corollary 4.4.} The elements of \(C\) of the form
\[
  \xi_1^{\alpha_1}\cdots \xi_{n-1}^{\alpha_{n-1}},
  \qquad 0\leq \alpha_i\leq n-i,
\]
form a basis of \(C\) as an \(A\)-module.

For the proof of the lemma and corollary we refer to [2], Exposé 4, pages 4.19, from which they are borrowed.

\sgaStableTarget{sga6.e06.sec.4.5}\subsection*{4.5}
Let \(S\) be quasi-compact, let \(E\) be an \(\mathcal O_S\)-module locally free of rank \(n\), let \(\pi=(p_1,\ldots,p_k)\) be a sequence of integers \(\geq 1\) with sum \(n\), and put \(X=D_\pi(E)\). The scheme \(X\) is equipped with canonical quotients \(L_i\). Put \(L_0=0\), \(L_k=f^*E\), where \(f:X\to S\) is the structural morphism. For \(1\leq i\leq k\), let \(F_i\) be the kernel of \(L_i\to L_{i-1}\). Let \(K^\bullet(X)\) be the Grothendieck group of locally free finite-type sheaves on \(X\), respectively of perfect complexes. Denote by \(F_j\) the class of \(F_j\), and by \(\lambda_i^{(j)}\) the class of \(\lambda^i(F_j)\). In the case where \(K^\bullet(X)\) is defined by perfect complexes, one gives the same evident meaning to the series \(\lambda_t(E)\) and \(\lambda_t(F_j)\), where \(E\) is the class of \(E\). Finally,
\[
  \sum_{j=1}^{k}F_j=E.
\]

\sgaStableTarget{sga6.e06.thm.4-6}\textbf{Theorem 4.6.} With the hypotheses and notations of \hyperlink{sga6.e06.sec.4.5}{4.5}, the \(K^\bullet(S)\)-algebra \(K^\bullet(X)\) is generated by the \(\lambda_i^{(j)}\), and the ideal of relations among the \(\lambda_i^{(j)}\)'s is generated by the relations deduced from
\[
  \lambda_t(E)=\prod_{j=1}^{k}\lambda_t(F_j).
\]
Equivalently, \(K^\bullet(X)\) is the \(K^\bullet(S)\)-\(\lambda\)-algebra generated by \(F_1,\ldots,F_k\), subject to the relations
\[
  \text{(i)}\quad \lambda^p(F_i)=0\quad(p>p_i),
  \qquad
  \text{(ii)}\quad \sum_i F_i=E.
\]

We first treat the case where all \(p_i\)'s are equal to \(1\). Then the theorem reads as follows.

\sgaStableTarget{sga6.e06.cor.4-7}\textbf{Corollary 4.7.} Let \(S\) be quasi-compact, let \(E\) be an \(\mathcal O_S\)-module locally free of rank \(n\), let \(X=D(E)\) be the flag bundle of \(E\), and let \(\xi_1,\ldots,\xi_n\) be the classes of the invertible sheaves \(F_1,\ldots,F_n\). Then \(K^\bullet(X)\) is the \(K^\bullet(S)\)-algebra generated by the \(\xi_i\)'s, subject to the relations
\[
  \sigma_i(\xi_1,\ldots,\xi_n)=\lambda^i(E).
\]

We argue by induction on \(n\). For \(n=1\) the assertion is trivial. Suppose it known for \(n-1\). Consider \(\mathbb P(E)\), and on \(\mathbb P(E)\) the exact sequence
\[
  0\longrightarrow F\longrightarrow g^*E\longrightarrow \mathcal O_{\mathbb P(E)}(1)\longrightarrow 0,
\]
where \(g\) is the structural morphism of \(\mathbb P(E)\). By \hyperlink{sga6.e06.lem.4-2}{Lemma 4.2}, \(X\) identifies with the usual flag bundle \(D(F)\) over \(\mathbb P(E)\), and the canonical sheaves correspond under this identification. The diagram
\[
\begin{tikzcd}[column sep=large,row sep=large]
D(E)=D(F) \arrow[r,"h"] \arrow[dr,"f"'] & \mathbb P(E) \arrow[d,"g"]\\
& S
\end{tikzcd}
\]
shows that \(K^\bullet(S)\to K^\bullet(X)\) factors through \(K^\bullet(\mathbb P(E))\). By the induction hypothesis applied to \(D(F)\), and by \hyperlink{sga6.e06.cor.4-4}{Corollary 4.4}, \(K^\bullet(X)\) has over \(K^\bullet(\mathbb P(E))\) a basis formed by
\[
  \xi_2^{\alpha_2}\cdots \xi_n^{\alpha_{n-1}},
  \qquad 0\leq \alpha_i\leq n-i.
\]
Moreover \(K^\bullet(\mathbb P(E))\) has over \(K^\bullet(S)\) the basis \(1,\xi_1,\ldots,\xi_1^{n-1}\). Hence \(K^\bullet(X)\) has over \(K^\bullet(S)\) the basis
\[
  \xi_1^{\alpha_1}\cdots \xi_{n-1}^{\alpha_{n-1}},
  \qquad 0\leq \alpha_i\leq n-i.
\]
In addition, the additivity of \(\lambda_t\) gives \(\prod_i(1+\xi_i t)=\lambda_t(E)\), so the displayed relations are satisfied. If \(C\) is the \(K^\bullet(S)\)-algebra generated by \(\xi_i\)'s and subjected to these relations, then there is a surjective homomorphism \(C\to K^\bullet(X)\). By \hyperlink{sga6.e06.cor.4-4}{Corollary 4.4} this homomorphism sends a basis of \(C\) to a basis of \(K^\bullet(X)\), hence it is an isomorphism.

For a general flag bundle of type \(\pi\), the usual flag bundle \(D(E)\) is obtained from \(D_\pi(E)\) by constructing successively the flag bundles of the \(F_j\)'s. It follows that \(K^\bullet(D(E))\), considered as a module over \(K^\bullet(X)\), has a basis formed by the elements appearing in 4.3.i), the \(\xi_i^{(j)}\)'s being the classes of the sheaves analogous to the \(F_j\)'s in the splitting of the inverse images of the \(F_j\)'s on \(D(E)\). Thus \(K^\bullet(X)\) is a subring of \(K^\bullet(D(E))\). Moreover, additivity gives
\[
  \lambda_i^{(j)}=\sigma_i(\xi_1^{(j)},\ldots,\xi_{p_j}^{(j)}),
\]
and therefore \(K^\bullet(X)\) contains the sub-\(K^\bullet(S)\)-algebra of \(K^\bullet(D(E))\) generated by the \(\lambda_i^{(j)}\)'s. If \(B\) is this subalgebra, then by 4.3.i) the same family of elements is a basis both for the \(B\)-module structure and for the \(K^\bullet(X)\)-module structure. Hence \(K^\bullet(X)=B\), proving generation. The ideal of relations among the \(\lambda_i^{(j)}\)'s is then given by 4.3.ii), taking into account the structure of \(K^\bullet(D(E))\) over \(K^\bullet(S)\) determined by \hyperlink{sga6.e06.cor.4-7}{4.7}, and the fact that the \(\xi_i^{(j)}\)'s are precisely the classes involved in the splitting of the inverse image of \(E\) on \(D(E)\). This proves the first form of 4.6.

For the second form, one observes that the elements \(F_i\) of the \(K^\bullet(S)\)-\(\lambda\)-algebra \(K^\bullet(X)\) satisfy (i) and (ii). If \(A\) is the \(K^\bullet(S)\)-\(\lambda\)-algebra generated by elements \(F_i'\) satisfying (i) and (ii), there is therefore a \(K^\bullet(S)\)-\(\lambda\)-homomorphism \(A\to K^\bullet(X)\) sending \(F_i'\) to \(F_i\). By \hyperlink{sga6.e06.thm.4-6}{4.6} this homomorphism is surjective. Since the \(\lambda\)-ideal generated by (i) and (ii) in \(K^\bullet(S)[F_1,\ldots,F_k]\) contains the usual ideal generated by the first form of the relations, the homomorphism is also injective. Thus it is an isomorphism.

\sgaStableTarget{sga6.e06.cor.4-8}\textbf{Corollary 4.8.} With the hypotheses and notations of \hyperlink{sga6.e06.sec.4.5}{4.5}, put
\[
  c_i^{(j)}=\gamma^i(F_j-p_j),
  \qquad 1\leq i\leq p_j.
\]
Then the \(K^\bullet(S)\)-algebra \(K^\bullet(X)\) is generated by the \(c_i^{(j)}\)'s, and the ideal of relations among them is generated by the relations deduced from
\[
  \gamma_t(E-n)=\prod_{j=1}^{k}\gamma_t(F_j-p_j).
\]
This is merely another formulation of 4.6. Indeed, the \(c_i^{(j)}\)'s may be expressed as functions of the \(\lambda_i^{(j)}\)'s and conversely. Moreover the relation \(\lambda_t(E)=\prod\lambda_t(F_j)\) is equivalent to \(\gamma_t(E-n)=\prod\gamma_t(F_j-p_j)\).

\sgaStableTarget{sga6.e06.sec.4.9}\subsection*{4.9}
Let \(S\) be a scheme. Denote by \(G_S^{n,r}\) the Grassmannian of type \((n,r-n)\) of \(\mathcal O_S^r\), with \(r\ge n\). For variable \(r\), the \(G_S^{n,r}\) form an inductive system, the morphism
\[
  G_S^{n,r}\longrightarrow G_S^{n,r+1}
\]
being defined by considering \(\mathcal O_{G_S^{n,r}}^r\) as the quotient of \(\mathcal O_{G_S^{n,r}}^{r+1}\) by the \((r+1)\)-st factor. In this system of morphisms the canonical quotients therefore correspond by inverse image. The inductive system so defined will be denoted \(G_S^{n,\infty}\).

We define its \(K^\bullet\) by the formula
\sgaStableTarget{sga6.e06.eq.4-9-1}\begin{equation}\tag{4.9.1}
  K^\bullet(G_S^{n,\infty})
  = \varprojlim_r K^\bullet(G_S^{n,r}) .
\end{equation}
Let \(F_1,F_2\) be the canonical sheaves on \(G_S^{n,r}\) \hyperlink{sga6.e06.sec.4.5}{(4.5)}, and \(F_1,F_2\) their classes in \(K^\bullet(G_S^{n,r})\). Put
\[
  c_i=\gamma^i(F_1-n).
\]
When \(r\) varies, the \(c_i\) correspond under the transition morphisms and consequently define elements, still denoted \(c_i\), of \(K^\bullet(G_S^{n,\infty})\).

\sgaStableTarget{sga6.e06.prop.4-10}\textbf{Proposition 4.10.} Let \(S\) be a quasi-compact scheme. Then, with the notation of \hyperlink{sga6.e06.sec.4.9}{4.9}, one has an isomorphism
\[
  K^\bullet(G_S^{n,\infty})\simeq K^\bullet(S)[[c_1,\ldots,c_n]] .
\]

By \hyperlink{sga6.e06.cor.4-8}{4.8}, \(K^\bullet(G_S^{n,r})\) is the \(K^\bullet(S)\)-algebra generated by
\[
  c_i=\gamma^i(F_1-n),\quad 1\le i\le n,
\]
and by
\[
  \gamma^j(F_2-r+n),\quad 1\le j\le r-n,
\]
subject to the relations deduced from
\[
  \gamma_t(F_1-n)\cdot\gamma_t(F_2-r+n)
  = \gamma_t(r-r)=1,
\]
or equivalently
\[
  \gamma_t(F_2-r+n)=1/\gamma_t(F_1-n).
\]
This relation makes it possible to express the \(\gamma^j(F_2-r+n)\) in terms of the \(c_i\), and gives a series of isobaric relations among the \(c_i\) (where \(c_i\) is assumed to have weight \(i\)) by writing that the coefficients of \(1/\gamma_t(F_1-n)\) vanish in degree \(>r-n\); moreover, these relations generate the ideal of relations among the \(c_i\). We may therefore write
\[
  K^\bullet(G_S^{n,r})=K^\bullet(S)[c_1,\ldots,c_n]/I_r,
\]
where \(I_r\) is a homogeneous ideal of \(K^\bullet(S)[c_1,\ldots,c_n]\), of degree \(>r-n\). It then follows from the \hyperlink{sga6.e06.eq.4-9-1}{definition (4.9.1)} that \(K^\bullet(G_S^{n,\infty})\) is the separated completion of \(K^\bullet(S)[c_1,\ldots,c_n]\) for the topology defined by the \(I_r\). This topology is finer than the topology defined by weight; we shall show that it is equivalent to it, which will imply the proposition.

Put \(A=K^\bullet(S)[c_1,\ldots,c_n]/I_r\), and let \(C\) be the \(A\)-algebra generated by generators \(x_1,\ldots,x_n\) subject to the relations
\[
  \sigma_j(x_1,\ldots,x_n)=c_j,
\]
where the \(\sigma_j\) denote the elementary symmetric functions. By \hyperlink{sga6.e06.cor.4-4}{4.4}, the canonical homomorphism \(A\to C\) is injective. We obtain
\[
  \gamma_t(F_1-n)
  =1+c_1t+\cdots+c_nt^n
  =\prod_{i=1}^n(1+x_it).
\]
Since \(1/\gamma_t(F_1-n)\) is a polynomial, so is \(1/(1+x_it)\), for every \(i\). Hence the \(x_i\) are nilpotent, and there exists an integer \(N\) such that every polynomial of degree \(>N\) in the \(x_i\), with coefficients in \(K^\bullet(S)\), is zero. Consequently every polynomial of \(A\) of weight \(>N\) is zero, i.e. every polynomial of \(K^\bullet(S)[c_1,\ldots,c_n]\) of weight \(>N\) belongs to \(I_r\). It follows that the two topologies under consideration are equivalent, completing the proof.

\sgaStableTarget{sga6.e06.sec.5}\section*{5. Applications to projective bundles; study of \(f_*\)}

\sgaStableTarget{sga6.e06.sec.5.1}\subsection*{5.1}
Let \(S\) be quasi-compact, let \(E\) be locally free of rank \(r+1\), and put \(X=\mathbb P(E)\), with structural morphism \(f:X\to S\). On \(X\) one has the canonical exact sequence
\[
  0\longrightarrow F\longrightarrow f^*E\longrightarrow \mathcal O_X(1)\longrightarrow 0.
\]
Let \(E,F,\xi\) denote the corresponding classes in \(K^\bullet\), and put \(x=\xi-1\).

\sgaStableTarget{sga6.e06.sec.5.2}\subsection*{5.2}
For a locally free sheaf \(G\) on \(X\), the higher direct images \(R^if_*G\) need not be locally free of finite type on \(S\). Thus the direct image \(f_*:K^\bullet(X)\to K^\bullet(S)\) is not defined here by alternating sums of higher direct images. \hyperlink{sga6.e06.thm.1-1}{Theorem 1.1} gives a substitute: every \(y\in K^\bullet(X)\) is uniquely expressible in the form
\[
  y=a_0+a_1\xi+\cdots+a_r\xi^r,
  \qquad a_i\in K^\bullet(S),
\]
and we put
\[
  f_*(y)=a_0.
\]
This additive homomorphism satisfies the projection formula. Moreover,
\[
  f_*(1)=1,
  \qquad
  f_*(\xi^n)=0\quad(-r\leq n<0),
  \qquad
  f_*(\xi^n)=s_n\quad(n\geq 0),
\]
where \(s_n\) is the \(n\)-th complete symmetric function in the roots of \(E\). The last relations follow by induction from the dependence relation and the Koszul relation.

\sgaStableTarget{sga6.e06.prop.5-3}\textbf{Proposition 5.3.} For every \(k\geq 0\),
\[
  \Fil^k K^\bullet(X)
  =
  \Fil^k K^\bullet(S)
  +\Fil^{k-1}K^\bullet(S)x+\cdots+\Fil^{k-r}K^\bullet(S)x^r,
\]
as a direct sum of \(K^\bullet(S)\)-modules, where the filtration is the canonical filtration of augmented \(\lambda\)-rings.

Here \(x\) has augmentation zero and
\[
  \gamma_t(x)=\lambda_{t/(1-t)}(1-\xi^{-1})=\frac{1}{1-xt},
\]
so that \(\gamma^i(x)=x^i\). The result follows from the computation of the filtration for free \(\lambda\)-rings in \hyperlink{sga6.e05.prop.4-12}{V, 4.12} and \hyperlink{sga6.e05.rem.4-13}{V, 4.13}, together with the relation \hyperlink{sga6.e06.eq.1-9-1}{(1.9.1)}. Directness follows from the independence of \(1,x,\ldots,x^r\).

\sgaStableTarget{sga6.e06.cor.5-4}\textbf{Corollary 5.4.} The filtration induced on \(K^\bullet(S)\) by \(K^\bullet(X)\) is the original filtration of \(K^\bullet(S)\).

\sgaStableTarget{sga6.e06.cor.5-5}\textbf{Corollary 5.5.} If \(X\) is a product over \(S\) of flag bundles of locally free sheaves, then
\[
  \Gr K^\bullet(S)\longrightarrow \Gr K^\bullet(X)
\]
is injective.

\sgaStableTarget{sga6.e06.cor.5-6}\textbf{Corollary 5.6.} With the notation of the projective bundle above, \(\Gr K^\bullet(X)\) is a free \(\Gr K^\bullet(S)\)-module with basis \(1,x,\ldots,x^r\), and
\[
  \Gr^k K^\bullet(X)
  =\Gr^kK^\bullet(S)\oplus\Gr^{k-1}K^\bullet(S)x\oplus\cdots\oplus\Gr^{k-r}K^\bullet(S)x^r.
\]
The elements \(1,x,\ldots,x^{r+1}\) satisfy
\[
  \sum_{i=0}^{r+1}(-1)^i c_i(E)x^{r+1-i}=0,
\]
where \(c_i(E)\) is the \(i\)-th Chern class of \(E\).

\sgaStableTarget{sga6.e06.cor.5-7}\textbf{Corollary 5.7.} For a flag bundle as in \hyperlink{sga6.e06.sec.4.5}{4.5}, the \(\Gr K^\bullet(S)\)-algebra \(\Gr K^\bullet(X)\) is generated by the classes \(\xi_i^{(\alpha)}\) of \(\xi_i^{(\alpha)}-1\), with relations generated by
\[
  c_t(E)=\prod_{i,\alpha}c_t(\xi_i^{(\alpha)}),
\]
where \(c_t\) denotes the total Chern class.

\sgaStableTarget{sga6.e06.cor.5-8}\textbf{Corollary 5.8.} For every \(k\),
\[
  f_*(\Fil^k K^\bullet(X))\subset \Fil^k K^\bullet(S).
\]
This follows immediately from \hyperlink{sga6.e06.prop.5-3}{Proposition 5.3} and the projection formula.

\sgaStableTarget{sga6.e06.prop.5-9}\textbf{Proposition 5.9.} With the notation of \hyperlink{sga6.e06.sec.5.1}{5.1}, for every \(y\in K^\bullet(X)\), one has
\[
  f_*(\lambda_{-1}(F)y)=0,
  \qquad
  f_*(\gamma^*(F)y)=0.
\]
In particular,
\[
  f_*(\lambda_{-1}(F))=1.
\]
If
\[
  \xi_t:1+K^\bullet(X)[[t]]^+\longrightarrow 1+K^\bullet(S)[[t]]^+
\]
is defined coefficientwise by
\[
  \xi_t\left(\sum a_i t^i\right)=\sum f_*(a_i)t^i,
\]
then, although it is not in general a homomorphism of groups, it satisfies
\[
  \xi_t(mf^*y)=\xi_t(m)y.
\]
Using \(E=F+\xi\) and \(\lambda_t(E)=\lambda_t(F)(1+\xi t)\), one obtains the asserted identities.

\sgaStableTarget{sga6.e06.prop.5-10}\textbf{Proposition 5.10.} Let \(z\in K^\bullet(X)\). The condition
\[
  z(1-\xi)=0
\]
is equivalent to the existence of an element \(y\in K^\bullet(S)\) such that
\[
  \text{(i)}\quad z=y\lambda_{-1}(F),
  \qquad
  \text{(ii)}\quad y\lambda_{-1}(E)=0.
\]
Sufficiency follows from
\[
  \lambda_{-1}(E)=\lambda_{-1}(F)\lambda_{-1}(\xi)=\lambda_{-1}(F)(1-\xi).
\]
Conversely, write \(z=\sum_{i=0}^r a_i\xi^i\). Combining the dependence relation among the powers of \(\xi\) with \(z(1-\xi)=0\), and using the independence of \(1,\ldots,\xi^r\), expresses \(z\) as a multiple of \(\lambda_{-1}(F)\). The top relation gives \(y\lambda_{-1}(E)=0\).

\sgaStableTarget{sga6.e06.prop.5-11}\textbf{Proposition 5.11.} Let \(z\in \Gr K^\bullet(X)\). The condition \(zx=0\) is equivalent to the existence of \(y\in \Gr K^\bullet(S)\) such that
\[
  \text{(i)}\quad z=y\,c_{\mathrm{top}}(F),
  \qquad
  \text{(ii)}\quad y\,c_{\mathrm{top}}(E)=0.
\]
The proof is the same as that of \hyperlink{sga6.e06.prop.5-10}{5.10}, after passing to associated graded rings and using 5.6.

\sgaStableTarget{sga6.e06.sec.6}\section*{6. Study of the filtration of \(K^\bullet(X)\) when \(X\) has an ample sheaf}

Let \(X\) be a scheme possessing an ample sheaf. Then the two definitions of \(K^\bullet(X)\), by locally free finite-type sheaves and by perfect complexes, coincide. Let
\[
  \varepsilon:K^\bullet(X)\longrightarrow H^0(X,\mathbb Z)
\]
be the augmentation.

\sgaStableTarget{sga6.e06.prop.6-1}\textbf{Proposition 6.1.} If \(X\) is quasi-compact and has an ample sheaf, then \(\Fil^1 K^\bullet(X)\) is a nilideal.

It is enough to show that \(E-\varepsilon(E)\) is nilpotent, where \(E\) is the class of a locally free sheaf. By decomposing \(X\) into finitely many disjoint opens and then passing to the flag bundle of \(E\), one reduces to the case where \(E\) is the class of an invertible sheaf. Let \(L\) be the class of an ample invertible sheaf. For \(n\) sufficiently large, \(E^\vee(n)\) is generated by global sections, giving a surjection \((E(-n))^m\to \mathcal O_X\) with locally free kernel. Thus in \(K^\bullet(X)\)
\[
  mEL^{-n}=1+F.
\]
The corresponding \(\gamma\)-relation shows that \(1-EL^{-n}\), and therefore \(L^n-E\), is nilpotent. Since \(E-1=(L^n-1)-(L^n-E)\), and since \(L^n-1\) is also nilpotent for \(n\) large, the assertion follows.

\sgaStableTarget{sga6.e06.prop.6-2}\textbf{Proposition 6.2.} Let \(R\) be an augmented sub-\(\lambda\)-algebra of \(K^\bullet(X)\), generated by the classes of finitely many sheaves. Then the \(\lambda\)-ring filtration of \(R\) is discrete and \(\Fil^1R\) is nilpotent. Moreover, \(R\) is a finite module over its image in \(H^0(X,\mathbb Z)\).

After reducing to constant ranks, \(R\) is generated by finitely many elements \(x_i=E_i-\varepsilon(E_i)\) and their \(\gamma\)-operations. \hyperlink{sga6.e06.prop.6-1}{Proposition 6.1} shows that all these elements are nilpotent. Hence the augmentation ideal of \(R\) is nilpotent; only finitely many monomials in the \(x_i\)'s and the \(\gamma^j(x_i)\)'s remain. This gives both discreteness and finiteness.

\sgaStableTarget{sga6.e06.cor.6-3}\textbf{Corollary 6.3.} Let \(A\) be a commutative ring. An element \(y\) of \(K^\bullet(X)\otimes_{\mathbb Z}A\) is invertible if and only if its augmentation is invertible. In particular, the class of every nowhere-zero locally free sheaf of finite type becomes invertible in \(K^\bullet(X)\otimes_{\mathbb Z}A\).

Indeed, the augmentation ideal is nil.

\sgaStableTarget{sga6.e06.cor.6-4}\textbf{Corollary 6.4.} Let \(X\) be quasi-compact with an ample sheaf, and let \(f:X'\to X\) be proper and perfect. Assume either that \(f\) is projective or that \(X\) is noetherian, so that the finiteness theorem applies. If there is an element \(x'\in K^\bullet(X')\otimes A\) such that the augmentation of \(f_*(x')\) is invertible, then
\[
  f^*:K^\bullet(X)\otimes A\longrightarrow K^\bullet(X')\otimes A
\]
is injective.
Indeed, by the projection formula,
\[
  f_*(f^*y\cdot x')=y\cdot f_*(x'),
\]
and \(f_*(x')\) is invertible by 6.3.

\sgaStableTarget{sga6.e06.sec.6.5}\subsection*{6.5. The codimension filtration}
Assume now that \(X\) is noetherian of dimension \(d\), and has an ample sheaf. Define a filtration \(K^\bullet(X)^i\) as follows: an element \(x\in K^\bullet(X)\) lies in \(K^\bullet(X)^i\) if, for every closed subscheme \(Y\subset X\), it can be represented by a perfect complex \(M^\bullet(Y)\) such that
\begin{enumerate}[label=\textup{\alph*)}]
\item \(\cl(M^\bullet)=x\);
\item \(\codim(\Supp H^j(M^\bullet)\cap Y,Y)\geq i+j\), and \(H^j(M^\bullet)=0\) for \(j\gg 0\).
\end{enumerate}
These \(K^\bullet(X)^i\)'s are additive subgroups.

\sgaStableTarget{sga6.e06.prop.6-6}\textbf{Proposition 6.6.} The filtration of \hyperlink{sga6.e06.sec.6.5}{6.5} is a ring filtration, and \(K^\bullet(X)^{d+1}=0\).

If \(x\) and \(x'\) are represented by perfect complexes \(M^\bullet\) and \(M'{}^\bullet\) satisfying the above codimension estimates, then \(xx'\) is represented by \(M^\bullet\otimes M'{}^\bullet\). The support of \(H^n(M^\bullet\otimes M'{}^\bullet)\) is contained in the union of the intersections of supports with total degree \(n\), and the codimension estimate follows from EGA \(0_{\mathrm{IV}}\), 14.2.2.3. Finally \(K^\bullet(X)^{d+1}=0\), because a support in codimension greater than \(d\) is empty.

\sgaStableTarget{sga6.e06.lem.6-7}\textbf{Lemma 6.7.} Let \(X=\Spec A\) be affine noetherian, let \(x_1,\ldots,x_n\) be finitely many points of \(X\), and let \(E\) be invertible. There exists a section \(s\in \Gamma(X,E)\) which is a basis of \(E_{x_i}\) for every \(i\).

The proof uses closed specializations \(x_i'\) with distinct maximal ideals \(\mathfrak m_i\), the isomorphisms \(M/\mathfrak m_iM\simeq A/\mathfrak m_i\), for \(M=\Gamma(X,E)\), and Nakayama's lemma.

\sgaStableTarget{sga6.e06.lem.6-8}\textbf{Lemma 6.8.} Let \(X\) be noetherian with an ample sheaf \(L\), let \(E\) be invertible, and let \(x_1,\ldots,x_n\) be points of \(X\). There exists an integer \(N\) such that, for every \(k\geq1\), \(E(kN)\) has a global section which is a basis at every \(x_i\).

Choose a power of the ample sheaf and a section whose non-vanishing locus is affine and contains all the \(x_i\)'s; apply \hyperlink{sga6.e06.lem.6-7}{Lemma 6.7} there, and extend after multiplying by a further power.

\sgaStableTarget{sga6.e06.thm.6-9}\textbf{Theorem 6.9.} Let \(X\) be a noetherian scheme of dimension \(d\) with an ample sheaf. Then
\[
  \Fil^{d+1}K^\bullet(X)=0,
\]
and therefore \((\Fil^1K^\bullet(X))^{d+1}=0\).

First, if \(E\) and \(E'\) are invertible, with classes still denoted \(E,E'\), then \(E-E'\in K^\bullet(X)^1\). For a closed \(Y\subset X\), choose one point in each irreducible component. \hyperlink{sga6.e06.lem.6-8}{Lemma 6.8} gives morphisms from suitable negative powers of the ample sheaf to \(E\) and \(E'\) which are isomorphisms at those points. The two-term complex built from these morphisms has class \(E-E'\), and is acyclic at the chosen generic points of \(Y\); hence its cohomology supports have codimension at least \(1\) in \(Y\).

Now take \(x\in \Fil^{d+1}K^\bullet(X)\). It is enough to treat products
\[
  \gamma^{i_1}(E_1-F_1)\cdots\gamma^{i_k}(E_k-F_k),
  \qquad \sum i_j\geq d+1,
\]
where the \(E_i,F_i\) are represented by locally free sheaves of the same constant rank. Pull back to the product of the flag bundles of these sheaves. The pullback is injective on \(K^\bullet\), and each class splits into sums of invertible classes. For invertible classes \(\xi,\eta\),
\[
  \gamma_t(\xi-\eta)
  =\frac{\gamma_t(\xi-1)}{\gamma_t(\eta-1)}
  =\frac{1+(\xi-1)t}{1+(\eta-1)t},
\]
so
\[
  \gamma^i(\xi-\eta)=(-1)^{i-1}(\xi-\eta)(\eta-1)^{i-1}\in K^\bullet(X')^i.
\]
Thus the pullback of \(x\) belongs to \(K^\bullet(X')^{d+1}\). Multiplication by a basis element of \(K^\bullet(X')\) over \(K^\bullet(X)\) of relative codimension \(r\) gives an element of \(K^\bullet(X')^{d+1+r}=0\), since \(\dim X'=d+r\). The basis property and injectivity then give \(x=0\).

\sgaStableTarget{sga6.e06.problem.6-10}\textbf{Problem 6.10.} Under the hypotheses of \hyperlink{sga6.e06.thm.6-9}{Theorem 6.9}, is it true that for every \(k\) one has
\[
  K^\bullet(X)^k=\Fil^kK^\bullet(X)?
\]

\sgaStableTarget{sga6.e06.bibliography-for-expose-vi}\subsection*{Bibliography for Exposé VI}
\begin{enumerate}[label={\textup{[\arabic*]}}]
\item H. Bass, \emph{K-Theory and stable algebra}, Publications Mathématiques de l'IHÉS, no. 22.
\item C. Chevalley, Séminaire 1958: \emph{Anneaux de Chow et applications}.
\item M. Hakim, \emph{Topos localement annelés et schémas relatifs}, thesis in preparation.
\item R. Hartshorne, \emph{Residues and duality}, Lecture Notes in Mathematics 20, Springer-Verlag.
\item J.-P. Serre, \emph{Le groupe de Grothendieck des schémas en groupes réductifs}, Publications Mathématiques de l'IHÉS, no. 34.
\end{enumerate}



\providecommand{\Kosz}{\mathrm K}
\providecommand{\Proj}{\operatorname{Proj}}
\providecommand{\Sym}{\operatorname{Sym}}
\providecommand{\gr}{\operatorname{gr}}
\providecommand{\Tor}{\operatorname{Tor}}
\providecommand{\Spec}{\operatorname{Spec}}
\providecommand{\Imf}{\operatorname{Im}}
\providecommand{\Ker}{\operatorname{Ker}}
\providecommand{\Id}{\operatorname{Id}}
\providecommand{\rad}{\operatorname{rad}}
\providecommand{\cO}{\mathcal O}
\providecommand{\cS}{\mathcal S}
\providecommand{\cJ}{\mathcal J}
\providecommand{\cN}{\mathcal N}
\providecommand{\bfP}{\mathbf P}
\providecommand{\bfA}{\mathbf A}


